\section{Implementation}
\label{sec:implementation}

We implement RAVEL in Python as an asynchronous control layer in front of
geo-distributed vLLM serving replicas~\cite{vllm}.
Our implementation targets vLLM~0.8.5.post1 and its frozen V0 scheduler
contract.
All control-plane operations that modify Router-owned queues, perform
replanning, or initiate dispatch are serialized by a single
\texttt{asyncio} scheduler lock.
New arrivals, mobility-hold expirations, engine feedback, and explicit
replanning requests therefore enter the same serialized scheduling path,
preventing concurrent passes from racing on tentative ownership or
prospective accounting.
Terminal callbacks may asynchronously release tracked service charges, but
any resulting replanning again enters this serialized path.

\paragraph{Serving profiles.}
RAVEL derives serving estimates from externally calibrated schema-v2 profiles
rather than hard-coded token rates.
Each profile carries an engine fingerprint and is bound to one serving cluster;
formal runs require a valid external profile for every cluster.
At runtime, RAVEL linearly interpolates Prefill and Decode measurements between
calibrated concurrency points and clamps queries outside the measured range to
the nearest endpoint.
The engine adapter additionally constructs a conservative Prefill envelope
from the maximum calibrated slope and intercept across concurrency levels.
For multi-stage completion requests, the Router may also load an offline
semantic output profile indexed only by request-visible workflow metadata,
such as $(\textit{stage\_id},\textit{stage\_num})$.
Neither profile exposes future arrivals or realized output length.

\paragraph{vLLM integration and local TBT control.}
RAVEL integrates with vLLM without modifying its source tree.
At process startup, a lightweight adapter feature-detects and wraps
\texttt{Scheduler.\_get\_priority} and
\texttt{Scheduler.\_schedule\_chunked\_prefill}.
If either required hook is unavailable, startup fails rather than silently
disabling RAVEL's local controls.
For LATENCY requests, the adapter encodes a Router-derived deadline-order key
together with the applicable TBT requirement in a reserved priority range;
non-LATENCY requests in the mixed path retain ordinary FCFS priority.
When LATENCY Decode is active, the adapter combines the tightest active TBT
requirement with the conservative Prefill envelope to derive a block-aligned
bound on the next Prefill quantum.
It temporarily reduces \texttt{max\_num\_batched\_tokens} for that
Chunked-Prefill scheduling iteration and restores the configured value in a
\texttt{finally} block.
The bound therefore affects only the current local scheduling decision and
does not introduce a separate Router-level placement policy.

\paragraph{Asynchronous handoff and failure handling.}
Engine handoff is implemented as an explicit asynchronous ownership
transition.
After Router-side admission and budget checks succeed, the target replica
first registers the request, after which RAVEL creates and tracks the
corresponding engine-facing posting task.
Successful replica registration together with successful posting-task
creation constitutes the handoff defined in
Section~\ref{sec:design:commitment}; only then does placement become final.

If replica registration is rejected, the request is restored to its
Router-owned queue together with its prospective state and remains eligible
for the ordinary mobility and replanning rules.
Failures while creating or executing an engine-facing posting task are
surfaced to the control plane rather than silently converted into new
Router-level placements.
Posting attempts carry attempt identifiers so that delayed callbacks from an
older attempt cannot modify ownership or accounting associated with a newer
one.

After successful handoff, replica-reported pending and running state continues
to inform subsequent Router quotes.
The corresponding active-service charge remains visible until terminal
completion.
At request termination, RAVEL releases that charge and requests another
serialized scheduling pass so that newly available capacity can affect
requests that remain Router-owned.

